%==============================================================================
% Dàn ý chi tiết - Chapter 1: Đặt vấn đề
%==============================================================================

\chapter{Đặt vấn đề}

%-----------------------------------------------------------------------------
\section{Giới thiệu}
%-----------------------------------------------------------------------------
\begin{itemize}
    \item Bối cảnh giáo dục hiện đại và nhu cầu tự động hóa trong kiểm tra đánh giá
    \item Tầm quan trọng của ngân hàng câu hỏi trong giảng dạy
    \item Giới thiệu tổng quan về đề tài nghiên cứu
    \item Đối tượng và phạm vi áp dụng của nghiên cứu
\end{itemize}

%-----------------------------------------------------------------------------
\section{Động lực nghiên cứu}
%-----------------------------------------------------------------------------
\begin{itemize}
    \item Khó khăn trong việc tạo câu hỏi chất lượng cao thủ công
    \item Hạn chế của các phương pháp sinh câu hỏi truyền thống
    \item Tiềm năng của LLM trong việc sinh nội dung tự động
    \item Nhu cầu phân loại câu hỏi theo Bloom's Taxonomy để đảm bảo độ khó và mục tiêu học tập
    \item Thách thức trong việc tạo câu hỏi với độ phức tạp cao (suy luận, phản biện)
    \item Vấn đề về chất lượng phương án nhiễu (distractors)
    \item Nhu cầu thực tế về ứng dụng có khả năng quét OCR từ tài liệu vật lý
\end{itemize}

%-----------------------------------------------------------------------------
\section{Mục tiêu nghiên cứu}
%-----------------------------------------------------------------------------
\begin{itemize}
    \item Mục tiêu tổng quát: Nghiên cứu và đề xuất phương pháp sinh câu hỏi trắc nghiệm tự động dựa trên mô hình ngôn ngữ lớn kết hợp hệ thống đa tác tử
    \item Các mục tiêu cụ thể:
    \begin{itemize}
        \item Nghiên cứu và áp dụng kiến trúc Multi-Agent (LangGraph) để sinh câu hỏi theo các cấp độ Bloom khác nhau
        \item Thiết kế bộ prompt cho từng cấp độ nhận thức (Retrieval, Inference/Synthesis, Critical Reasoning)
        \item Xây dựng quy trình OCR để chuyển đổi văn bản thành hình ảnh và nhận dạng lại
        \item Đề xuất khung đánh giá đa chiều (Micro, Macro, Format, OCR)
        \item Xây dựng ứng dụng minh họa Openlet
    \end{itemize}
    \item Đóng góp chính của luận văn
\end{itemize}

%-----------------------------------------------------------------------------
\section{Phạm vi nghiên cứu}
%-----------------------------------------------------------------------------
\begin{itemize}
    \item Phạm vi về dữ liệu: Sử dụng tập dữ liệu RACE (Reading Comprehension)
    \item Phạm vi về mô hình: Các LLM phổ biến (GPT-4, Claude, LLaMA, etc.)
    \item Phạm vi về cấp độ Bloom: 3 cấp độ (Retrieval, Inference/Synthesis, Critical Reasoning)
    \item Phạm vi về phương pháp: So sánh Single-Prompt vs Multi-Agent
    \item Các giới hạn và ràng buộc của nghiên cứu
    \item Những vấn đề không thuộc phạm vi nghiên cứu
\end{itemize}

%-----------------------------------------------------------------------------
\section{Bố cục luận văn}
%-----------------------------------------------------------------------------
\begin{itemize}
    \item Mô tả tóm tắt nội dung từng chương
    \begin{itemize}
        \item Chương 1: Đặt vấn đề
        \item Chương 2: Cơ sở lý thuyết
        \item Chương 3: Phương pháp đề xuất
        \item Chương 4: Thực nghiệm và đánh giá
        \item Chương 5: Kết luận và hướng phát triển
    \end{itemize}
    \item Sơ đồ tổ chức luận văn
\end{itemize}
